\QHeading{第9题　脱式计算。}
\par\noindent


\CalcBlock[方法一：简便计算]{1}{\begin{aligned}[t]&\phantom{={}}\color{black}6.5-2.18+2.09\\&\color{answerred}=6.5-(2.18-2.09)\\&\color{answerred}=6.5-0.09\\&\color{answerred}=6.41\end{aligned}}
\par

\CalcBlock[方法一]{2}{\begin{aligned}[t]&\phantom{={}}\color{black}5.23-(1.26+3.27)\\&\color{answerred}=5.23-4.53\\&\color{answerred}=0.7\end{aligned}}
\MethodBlock{方法二：去括号后计算}{加号后的括号去掉不变号，减号后的括号去掉要改变各项的加减号。}{\begin{aligned}[t]&\phantom{={}}5.23-(1.26+3.27)\\
&=5.23-1.26-3.27\\
&=3.97-3.27\\
&=0.7\end{aligned}}

\CalcBlock{3}{\begin{aligned}[t]&\phantom{={}}\color{black}14-3.8+6.9\\&\color{answerred}=10.2+6.9\\&\color{answerred}=17.1\end{aligned}}
\CalcBlock[方法一]{4}{\begin{aligned}[t]&\phantom{={}}\color{black}32.25-(9.6-4.98)\\&\color{answerred}=32.25-4.62\\&\color{answerred}=27.63\end{aligned}}
\MethodBlock{方法二：去括号后计算}{加号后的括号去掉不变号，减号后的括号去掉要改变各项的加减号。}{\begin{aligned}[t]&\phantom{={}}32.25-(9.6-4.98)\\
&=32.25-9.6+4.98\\
&=22.65+4.98\\
&=27.63\end{aligned}}

\CalcBlock{5}{\begin{aligned}[t]&\phantom{={}}\color{black}18-6.4+0.89\\&\color{answerred}=11.6+0.89\\&\color{answerred}=12.49\end{aligned}}
\CalcBlock[方法一：简便计算]{6}{\begin{aligned}[t]&\phantom{={}}\color{black}4.25+7.2-3.6\\&\color{answerred}=4.25+(7.2-3.6)\\&\color{answerred}=4.25+3.6\\&\color{answerred}=7.85\end{aligned}}
\par\noindent

\QHeading{第10题　解决问题。}

\Problem{① 水泥厂计划全月生产8350吨水泥，实际上半月生产了4217.86吨，下半月生产了4479.57吨。水泥厂这个月实际比计划多生产了多少吨水泥？}
\par


\MethodTitle{方法一}
\Answer{{\Formula{\begin{aligned}[t]&\phantom{={}}4217.86+4479.57-8350\\&=8697.43-8350\\&=347.43\text{（吨）}\end{aligned}}}\par 答：多生产347.43吨。}

\MethodTitle{方法二：先求下半月还需完成的计划量}
\Answer{总计划减去上半月实际产量，得到下半月原本还需完成的数量，再与下半月实际产量比较。}
\Answer{\Formula{\begin{aligned}[t]&\phantom{={}}4479.57-(8350-4217.86)\\
&=4479.57-4132.14\\
&=347.43\text{（吨）}\end{aligned}}}
\Answer{答：实际比计划多生产347.43吨水泥。}
\Problem{② 把一根长8.45米的钢材锯成三段，第一段长1.88米，第二段长2.25米，第三段长多少米？}

\MethodTitle{方法一：简便计算}
\Answer{{\Formula{\begin{aligned}[t]&\phantom{={}}8.45-1.88-2.25\\&=8.45-2.25-1.88\\&=6.2-1.88\\&=4.32\text{（米）}\end{aligned}}}\par 答：第三段长4.32米。}

\MethodTitle{方法二：逐段减去}
\Answer{先减第一段，再减第二段。}
\Answer{\Formula{\begin{aligned}[t]&\phantom{={}}8.45-1.88-2.25\\
&=6.57-2.25\\
&=4.32\text{（米）}\end{aligned}}}
\Answer{答：第三段长4.32米。}
\Problem{③ 修路队第一天修路6.85千米，第二天比第一天多修1.3千米，第三天修的比前两天修的总和少6.97千米。第三天修了多少千米？}

\MethodTitle{方法一}
\Answer{{\Formula{\begin{aligned}[t]&\phantom{={}}6.85+(6.85+1.3)-6.97\\&=6.85+8.15-6.97\\&=15-6.97\\&=8.03\text{（千米）}\end{aligned}}}\par 答：第三天修了8.03千米。}

\MethodTitle{方法二：用两份第一天的长度加差额}
\Answer{第二天比第一天多1.3千米，所以前两天的总长是第一天的2倍再加1.3千米。}
\Answer{\Formula{\begin{aligned}[t]&\phantom{={}}6.85\times 2+1.3-6.97\\
&=13.7+1.3-6.97\\
&=15-6.97\\
&=8.03\text{（千米）}\end{aligned}}}
\Answer{答：第三天修了8.03千米。}
